% ============================================================
% Phần SH-A: Nền tảng toán học đầy đủ của Spherical Harmonics
% 9 slide, từ phương trình Laplace đến hàm điều hòa cầu thực
% File này KHÔNG tự đứng độc lập -- được \input trong main.tex
% ============================================================

% --- 1_1_laplace_equation.png ---
\begin{frame}{Xuất phát điểm: Phương trình Laplace}
\begin{columns}[T]
\begin{column}{0.52\textwidth}
\footnotesize
\begin{itemize}
    \item Hàm điều hòa (harmonic function) $\Phi(x,y,z)$ là nghiệm của phương trình Laplace trong không gian 3 chiều:
    $$\nabla^2\Phi=\frac{\partial^2\Phi}{\partial x^2}+\frac{\partial^2\Phi}{\partial y^2}+\frac{\partial^2\Phi}{\partial z^2}=0$$
    \item Chuyển sang toạ độ cầu $(r,\theta,\phi)$ với $x=r\sin\theta\cos\phi,\ y=r\sin\theta\sin\phi,\ z=r\cos\theta$, toán tử Laplace trở thành:
    $$\nabla^2\Phi=\dfrac{1}{r^2}\dfrac{\partial}{\partial r}\!\left(r^2\dfrac{\partial\Phi}{\partial r}\right)+\dfrac{1}{r^2\sin\theta}\dfrac{\partial}{\partial\theta}\!\left(\sin\theta\dfrac{\partial\Phi}{\partial\theta}\right)+\dfrac{1}{r^2\sin^2\theta}\dfrac{\partial^2\Phi}{\partial\phi^2}=0$$
    \item Phần góc $(\theta,\phi)$ của nghiệm chính là hàm điều hòa cầu $Y_l^m$ -- lý do SH là ``cơ sở tự nhiên'' cho hàm số trên mặt cầu, tương tự chuỗi Fourier là cơ sở tự nhiên trên đường tròn.
    \item Kiểm chứng số học (sai phân hữu hạn) trên lưới 2D lát cắt $z=\text{const}$: $\Phi=x^2-y^2$ cho $\nabla^2\Phi\approx0$, trong khi $\Phi=x^2+y^2$ cho $\nabla^2\Phi=4\neq0$.
    \item Đây là điểm khởi đầu của toàn bộ chuỗi suy diễn: Laplace $\to$ tách biến $\to$ Legendre $\to$ $Y_l^m$.
\end{itemize}
\end{column}
\begin{column}{0.44\textwidth}
\centering
\includegraphics[width=0.9\linewidth]{1_1_laplace_equation.png}
\captionof{figure}{\footnotesize Kiểm chứng số học $\nabla^2\Phi=0$ (hàm điều hòa) so với $\nabla^2\Phi\neq0$ trên lưới 2D.}
\end{column}
\end{columns}
\end{frame}

% --- 1_2_spherical_coords.png ---
\begin{frame}{Hệ toạ độ cầu $(\theta,\phi)$}
\begin{columns}[T]
\begin{column}{0.52\textwidth}
\begin{itemize}
    \item Toạ độ cầu dùng góc thiên đỉnh (polar angle) $\theta\in[0,\pi]$ và góc phương vị (azimuthal angle) $\phi\in[0,2\pi)$.
    \item Quan hệ với toạ độ Descartes trên mặt cầu bán kính $r$:
    $$
    \begin{aligned}
    x&=r\sin\theta\cos\phi\\
    y&=r\sin\theta\sin\phi\\
    z&=r\cos\theta
    \end{aligned}
    $$
    \item Với mặt cầu đơn vị ($r=1$): $(x,y,z)=(\sin\theta\cos\phi,\ \sin\theta\sin\phi,\ \cos\theta)$.
    \item Các hệ số $\sin\theta,\ \cos\theta,\ \sin^2\theta$ xuất hiện trực tiếp trong toán tử Laplace cầu -- đây là nguồn gốc của trọng số đo (measure) $d\Omega=\sin\theta\,d\theta\,d\phi$ dùng trong tích phân trực chuẩn về sau.
    \item Lưới toạ độ cầu $(\theta,\phi)$ được dựng bằng \texttt{meshgrid} và chiếu lên mặt cầu đơn vị để trực quan hoá.
\end{itemize}
\end{column}
\begin{column}{0.44\textwidth}
\centering
\includegraphics[width=0.9\linewidth]{1_2_spherical_coords.png}
\captionof{figure}{\footnotesize Lưới toạ độ cầu $(\theta,\phi)$ trên mặt cầu đơn vị.}
\end{column}
\end{columns}
\end{frame}

% --- 1_2b_radial_solution.png ---
\begin{frame}{Tách biến: Nghiệm phần bán kính $R(r)$}
\begin{columns}[T]
\begin{column}{0.52\textwidth}
\footnotesize
\begin{itemize}
    \item Giả sử nghiệm tách được: $\Phi(r,\theta,\phi)=R(r)\,\Theta(\theta)\,\Phi_m(\phi)$. Nhân hai vế phương trình Laplace với $r^2/(R\Theta\Phi_m)$, ta tách thành phần chỉ phụ thuộc $r$ và phần chỉ phụ thuộc $(\theta,\phi)$:
    $$\underbrace{\frac{1}{R}\frac{d}{dr}\!\left(r^2\frac{dR}{dr}\right)}_{\text{chỉ phụ thuộc }r}+\underbrace{\frac{1}{\Theta\sin\theta}\frac{d}{d\theta}\!\left(\sin\theta\frac{d\Theta}{d\theta}\right)+\frac{1}{\Phi_m\sin^2\theta}\frac{d^2\Phi_m}{d\phi^2}}_{\text{chỉ phụ thuộc }\theta,\phi}=0$$
    \item Vì hai nhóm phụ thuộc biến độc lập nhau, mỗi nhóm phải bằng một hằng số đối nhau. Đặt hằng số phân ly là $l(l+1)$:
    $$\frac{d}{dr}\!\left(r^2\frac{dR}{dr}\right)=l(l+1)R$$
    \item Nghiệm tổng quát của phương trình Euler-Cauchy này là hai nhánh độc lập:
    $$R(r)=A\,r^{l}+\dfrac{B}{r^{l+1}}$$
    \item Nhánh $r^l$ hữu hạn tại gốc toạ độ $r=0$; nhánh $r^{-l-1}$ hữu hạn tại vô cực -- tuỳ điều kiện biên vật lý mà chọn nhánh phù hợp.
\end{itemize}
\end{column}
\begin{column}{0.44\textwidth}
\centering
\includegraphics[width=0.9\linewidth]{1_2b_radial_solution.png}
\captionof{figure}{\footnotesize Hai nhánh nghiệm bán kính $R(r)=Ar^l$ và $R(r)=B/r^{l+1}$ với $l=0,1,2,3$.}
\end{column}
\end{columns}
\end{frame}

% --- 1_3_azimuthal_solution.png ---
\begin{frame}{Tách biến: Nghiệm phần phương vị $\Phi_m(\phi)$}
\begin{columns}[T]
\begin{column}{0.52\textwidth}
\begin{itemize}
    \item Từ phần góc còn lại, đặt hằng số tách thứ hai:
    $$-\dfrac{1}{\Phi_m}\dfrac{d^2\Phi_m}{d\phi^2}=m^2$$
    \item Phương trình vi phân thu được: $\dfrac{d^2\Phi_m}{d\phi^2}=-m^2\Phi_m$, có nghiệm tổng quát dạng sóng phức:
    $$\Phi_m(\phi)=e^{im\phi}$$
    \item Điều kiện tuần hoàn vật lý bắt buộc $\Phi_m(\phi+2\pi)=\Phi_m(\phi)$, tức $e^{im\cdot2\pi}=1$, dẫn đến điều kiện lượng tử hoá:
    $$m\in\mathbb{Z}$$
    \item Phần thực $\cos(m\phi)$ và quỹ đạo phức $e^{im\phi}=\cos(m\phi)+i\sin(m\phi)$ trên vòng tròn đơn vị được minh hoạ cho các giá trị nguyên $m=0,1,2,3$.
    \item Đây chính là "phiên bản cầu" của chuỗi Fourier theo góc phương vị -- bước trung gian trước khi ghép với nghiệm Legendre theo $\theta$.
\end{itemize}
\end{column}
\begin{column}{0.44\textwidth}
\centering
\includegraphics[width=0.9\linewidth]{1_3_azimuthal_solution.png}
\captionof{figure}{\footnotesize $\cos(m\phi)$ và quỹ đạo $e^{im\phi}$ trên vòng tròn đơn vị, $m=0,1,2,3$.}
\end{column}
\end{columns}
\end{frame}

% --- 1_4_legendre_functions.png ---
\begin{frame}{Phương trình Legendre liên kết}
\begin{columns}[T]
\begin{column}{0.52\textwidth}
\footnotesize
\begin{itemize}
    \item Đặt $x=\cos\theta$, phần góc $\theta$ còn lại trở thành \textbf{phương trình Legendre liên kết} (associated Legendre equation):
    $$\frac{d}{dx}\!\left[(1-x^2)\frac{d\Theta}{dx}\right]+\left[l(l+1)-\frac{m^2}{1-x^2}\right]\Theta=0$$
    \item Nghiệm hữu hạn tại hai cực $x=\pm1$ (cực Bắc/Nam của mặt cầu) chỉ tồn tại khi $l=0,1,2,\dots$ và $m=-l,\dots,l$ nguyên -- đây chính là nguồn gốc của điều kiện lượng tử hoá bậc $l$.
    \item Nghiệm là \textbf{hàm Legendre liên kết}:
    $$P_l^m(x)=(-1)^m(1-x^2)^{m/2}\frac{d^m}{dx^m}P_l(x)$$
    với đa thức Legendre thường (trường hợp $m=0$) cho bởi công thức Rodrigues:
    $$P_l(x)=\frac{1}{2^l\,l!}\frac{d^l}{dx^l}(x^2-1)^l$$
    \item Đồ thị minh hoạ: đa thức Legendre $P_l(x)$ với $l=0,\dots,5$, và hàm Legendre liên kết $P_3^m(x)$ với $m=0,\dots,3$.
\end{itemize}
\end{column}
\begin{column}{0.44\textwidth}
\centering
\includegraphics[width=0.9\linewidth]{1_4_legendre_functions.png}
\captionof{figure}{\footnotesize $P_l(x)$ với $l=0..5$ và $P_3^m(x)$ với $m=0..3$.}
\end{column}
\end{columns}
\end{frame}

% --- 1_5_spherical_harmonics_lobes.png ---
\begin{frame}{Ghép lại: Định nghĩa đầy đủ $Y_l^m$}
\begin{columns}[T]
\begin{column}{0.52\textwidth}
\footnotesize
\begin{itemize}
    \item Ghép nghiệm phần Legendre $P_l^m(\cos\theta)$ với nghiệm phương vị $e^{im\phi}$, cùng hằng số chuẩn hoá để đảm bảo trực chuẩn, ta có định nghĩa đầy đủ của hàm điều hòa cầu (spherical harmonics):
    $$\boxed{Y_l^m(\theta,\phi)=\sqrt{\dfrac{2l+1}{4\pi}\dfrac{(l-m)!}{(l+m)!}}\;P_l^m(\cos\theta)\,e^{im\phi}}$$
    \item $l\ge0$ gọi là \textbf{bậc} (degree), $m\in\{-l,\dots,l\}$ gọi là \textbf{bậc con} (order) -- ứng với mỗi $l$ có $2l+1$ hàm $Y_l^m$ độc lập.
    \item Hình dạng "thuỳ" (lobe) của $|Y_l^m|$ trên mặt cầu phụ thuộc vào cặp $(l,m)$: $l$ càng lớn, số thuỳ càng nhiều (dao động không gian càng nhanh); $m$ quyết định phân bố thuỳ theo phương vị.
    \item Trực quan hoá: bán kính tại mỗi hướng $=|Y_l^m(\theta,\phi)|$, màu sắc thể hiện dấu/giá trị của phần thực $\mathrm{Re}(Y_l^m)$.
\end{itemize}
\end{column}
\begin{column}{0.44\textwidth}
\centering
\includegraphics[width=0.9\linewidth]{1_5_spherical_harmonics_lobes.png}
\captionof{figure}{\footnotesize Hình dạng thuỳ của $\mathrm{Re}(Y_l^m)$ trên mặt cầu cho các cặp $(l,m)$ tiêu biểu.}
\end{column}
\end{columns}
\end{frame}

% --- 1_6_eigenvalue_check.png ---
\begin{frame}{$Y_l^m$ là hàm riêng của Laplace-Beltrami}
\begin{columns}[T]
\begin{column}{0.52\textwidth}
\begin{itemize}
    \item Hàm điều hòa cầu $Y_l^m$ chính là \textbf{hàm riêng} (eigenfunction) của toán tử Laplace-Beltrami $\nabla^2_{S^2}$ trên mặt cầu đơn vị $S^2$ (phần góc của toán tử Laplace 3D), với \textbf{trị riêng} (eigenvalue) $-l(l+1)$:
    $$\nabla^2_{S^2}Y_l^m=-l(l+1)\,Y_l^m$$
    \item Đây là hệ quả trực tiếp của quá trình tách biến: $l(l+1)$ chính là hằng số phân ly đã đặt ra ở bước tách phương trình Laplace 3D thành phần bán kính và phần góc.
    \item Kiểm chứng số học bằng sai phân hữu hạn của $\nabla^2_{S^2}$ trên lưới $(\theta,\phi)$:
    $$\nabla^2_{S^2}Y=\frac{1}{\sin\theta}\frac{\partial}{\partial\theta}\!\left(\sin\theta\frac{\partial Y}{\partial\theta}\right)+\frac{1}{\sin^2\theta}\frac{\partial^2Y}{\partial\phi^2}$$
    tính bằng \texttt{np.gradient} hai lần, sau đó so sánh tỉ số $\nabla^2_{S^2}Y/Y$ với giá trị lý thuyết $-l(l+1)$.
    \item Kết quả số khớp với lý thuyết trong sai số rời rạc hoá -- xác nhận $Y_l^m$ đúng là nghiệm riêng, không phải xấp xỉ.
\end{itemize}
\end{column}
\begin{column}{0.44\textwidth}
\centering
\includegraphics[width=0.9\linewidth]{1_6_eigenvalue_check.png}
\captionof{figure}{\footnotesize Kiểm chứng số học trị riêng $-l(l+1)$ bằng sai phân hữu hạn của $\nabla^2_{S^2}Y_l^m$.}
\end{column}
\end{columns}
\end{frame}

% --- 1_7_orthonormality_matrix.png ---
\begin{frame}{Tính trực chuẩn của hệ $\{Y_l^m\}$}
\begin{columns}[T]
\begin{column}{0.52\textwidth}
\scriptsize
\begin{itemize}
    \item Hệ $\{Y_l^m\}$ thoả \textbf{trực chuẩn} trên $S^2$ với đo $d\Omega{=}\sin\theta\,d\theta\,d\phi$:
    $\int_{S^2}Y_l^m(Y_{l'}^{m'})^{*}d\Omega{=}\delta_{ll'}\delta_{mm'}$.
    \item Hệ quả Sturm-Liouville (phương trình Legendre, phần $\theta$) kết hợp trực giao $e^{im\phi}$ trên chu kỳ $2\pi$ (phần $\phi$).
    \item Kiểm chứng số: ma trận Gram $M_{ij}{=}\sum Y_i(Y_j)^{*}\sin\theta\,\Delta\theta\Delta\phi \approx \delta_{ij}$ cho $l{=}0,1,2$.
    \item Nền tảng dùng $\{Y_l^m\}$ làm cơ sở Fourier trên mặt cầu — hệ số tính bằng phép chiếu, không chồng lấn giữa các bậc.
\end{itemize}
\end{column}
\begin{column}{0.44\textwidth}
\centering
\includegraphics[width=0.9\linewidth]{1_7_orthonormality_matrix.png}
\captionof{figure}{\footnotesize Ma trận Gram xấp xỉ ma trận đơn vị.}
\end{column}
\end{columns}
\end{frame}

% --- 2_1_real_sh_construction.png ---
\begin{frame}{Từ SH phức sang SH thực (Real SH)}
\begin{columns}[T]
\begin{column}{0.55\textwidth}
\scriptsize
    Đồ hoạ máy tính (3DGS) dùng số thực, nhưng $Y_l^m$ gốc chứa $e^{im\phi}$ là số phức $\Rightarrow$ tổ hợp tuyến tính $Y_l^{\pm m}$ để có cơ sở \textbf{thực} tương đương, giữ trực chuẩn:
    \[
    Y_{lm}^{\text{real}}=
    \begin{cases}
    \frac{i}{\sqrt{2}}\left(Y_l^{m}-(-1)^m Y_l^{-m}\right) & m<0\\
    Y_l^0 & m=0\\
    \frac{1}{\sqrt{2}}\left(Y_l^{-m}+(-1)^m Y_l^{m}\right) & m>0
    \end{cases}
    \]
    \begin{itemize}
    \item Với $m\neq0$: tổ hợp $Y_l^m$ và liên hợp $Y_l^{-m}=(-1)^m(Y_l^m)^*$ triệt tiêu phần ảo.
    \item \textbf{Lý do dùng SH thực trong 3DGS}: tránh số phức toàn pipeline, giảm nửa số biến lưu, khớp trực tiếp RGB trong \texttt{computeColorFromSH}.
    \end{itemize}
\end{column}
\begin{column}{0.45\textwidth}
\centering
\includegraphics[width=0.95\linewidth]{2_1_real_sh_construction.png}
\captionof{figure}{\footnotesize Dựng $Y_{1,m}^{\text{real}}$, $m{=}-1,0,1$.}
\end{column}
\end{columns}
\end{frame}
